%% kapa: exam=1 round=35 year=2024 period=2 form=A subject=appraisal_law subject_ko=감정평가관계법규 seq=24 number=24
%% source: Q-Net 감정평가사 기출문제 (cst00309) artlSeq=5213229 file=raw/1008/5213229/2교시 문제지 원본.pdf
%% conversion: kordoc 4.24.51 -> tools/gen1.py (auto); review=proofread-sonnet; flags=none
\begin{question}{24}
건축법령상 용어의 정의에 관한 설명으로 옳지 않은 것은?
\begin{choices}
  \item 기존 건축물의 전부를 해체하고 그 대지에 종전과 같은 규모의 범위에서 건축물을 다시 축조하는 것은 “개축”에 해당한다.
  \item “재축”에 해당하려면 연면적 합계는 종전 규모 이하로 하여야 한다.
  \item “이전”이란 건축물의 주요구조부를 해체하지 아니하고 같은 대지의 다른 위치로 옮기는 것을 말한다.
  \item 16층 이상인 건축물은 그 용도에 관계없이 “다중이용 건축물”이다.
  \item 기둥과 기둥 사이의 거리가 15미터 이상인 건축물은 “특수구조 건축물”이다.
\end{choices}
\end{question}
\answer{5}
